\subsection{Requisiti funzionali}
\label{subsec:requisiti_funzionali}
I requisiti funzionali sono raggruppati usando una divisione ad alto livello del sistema in due componenti, ovvero 
\begin{enumerate}
    \item Logica di presentazione: si occupa della presentazione delle informazioni agli utenti.
    \item Logica di business: si occupa dell'elaborazione delle informazioni mostrate agli utenti e ottenute dagli utenti.
\end{enumerate}

\subsubsection{Logica di presentazione}

\paragraph{Obbligatori}

\begin{freq}
    {RFO-1}
    {Il sistema deve poter visualizzare gli elementi necessari per avviare un'analisi}
    \label{rf:RFO-1}%

    \subreq{RFO-1.1}{\hyperref[uc:UC-1]{UC-1}}{Il sistema deve visualizzare un pulsante che permetta di iniziare l'analisi delle immagini senza alt}
    \subreq{RFO-1.2}{\hyperref[uc:UC-6]{UC-6}}{Il sistema deve visualizzare un pulsante che permetta di iniziare l'analisi dei video senza sottotitoli}
\end{freq}

\begin{freq}
    {RFO-2}
    {Il sistema deve poter visualizzare il risultato dell'analisi delle immagini}
    \label{rf:RFO-2}%

    \subreq{RFO-2.1}{\hyperref[uc:UC-2]{UC-2}}{Se non sono presenti immagini senza alt il sistema deve mostrare un messaggio che lo indica}
    \subreq{RFO-2.2}{\hyperref[uc:UC-2]{UC-2}}{Se sono presenti immagini senza alt il sistema deve mostrare la prima pagina di risultati dell'analisi seguendo \hyperref[rf:RFO-4]{RFO-4}}
\end{freq}

\begin{freq}
    {RFO-3}
    {Il sistema deve poter visualizzare il risultato dell'analisi dei video}
    \label{rf:RFO-3}%

    \subreq{RFO-3.1}{\hyperref[uc:UC-2]{UC-2}}{Se non sono presenti video senza sottotitoli il sistema deve mostrare un messaggio che lo indica}
    \subreq{RFO-3.2}{\hyperref[uc:UC-2]{UC-2}}{Se sono presenti video senza sottotitoli il sistema deve mostrare la prima pagina di risultati dell'analisi seguendo \hyperref[rf:RFO-4]{RFO-4}}
\end{freq}

\begin{freq}
    {RFO-4}
    {Il sistema deve poter visualizzare una pagina dei risultati di un'analisi di immagini o video}
    \label{rf:RFO-4}%

    \subreq{RFO-4.1}{\hyperref[uc:UC-2]{UC-2}}{Il sistema deve visualizzare la pagina dei risultati sottoforma di una lista scorrevole di elementi, seguendo \hyperref[rf:RFO-5]{RFO-5}}
    \subreq{RFO-4.2}{\hyperref[uc:UC-2]{UC-2}}{Il sistema deve visualizzare il numero totale delle pagine di risultati}
    \subreq{RFO-4.3}{\hyperref[uc:UC-2]{UC-2}}{Il sistema deve visualizzare il numero della pagina visualizzata}
    \subreq{RFO-4.4}{\hyperref[uc:UC-2]{UC-2}}{Il sistema deve visualizzare i pulsanti necessari per cambiare pagina di risultati}
\end{freq}

\begin{freq}
    {RFO-5}
    {Il sistema deve poter visualizzare un risultato dell'analisi}
    \label{rf:RFO-5}%

    \subreq{RFO-5.1}{\hyperref[uc:UC-2]{UC-2}}{Nel caso di un'immagine, il sistema deve permettere la visualizzazione dell'immagine analizzata}
    \subreq{RFO-5.2}{\hyperref[uc:UC-2]{UC-2}}{Nel caso di un video, il sistema deve permettere la visualizzazione, se presente, dell'immagine assegnata come poster al video, altrimenti un'immagine di default}
    \subreq{RFO-5.3}{\hyperref[uc:UC-2]{UC-2}}{Il sistema deve mostrare il valore dell'attributo \texttt{src} dell'immagine o del video}
    \subreq{RFO-5.4}{\hyperref[uc:UC-2]{UC-2}, \hyperref[uc:UC-4]{UC-4}}{Nel caso di un'immagine, il sistema deve mostrare l'\texttt{alt} suggerito}
    \subreq{RFO-5.5}{\hyperref[uc:UC-2]{UC-2}, \hyperref[uc:UC-7]{UC-7}}{Nel caso di un video, il sistema deve mostrare i sottotitoli suggeriti}

\end{freq}

\begin{freq}
    {RFO-6}
    {In caso di errore durante l'analisi, il sistema deve visualizzare un messaggio informativo}
    \label{rf:RFO-6}%

    \subreq{RFO-6.1}{\hyperref[uc:UC-5]{UC-5}, \hyperref[uc:UC-8]{UC-8}}{Il sistema deve mostrare un messaggio informativo in caso di errore nell'interazione con l'LLM}
\end{freq}

\paragraph{Facoltativi}

\begin{freq}
    {RFF-1}
    {Il sistema deve poter visualizzare gli elementi necessari per avviare un'analisi dell'ordine della sequenza di tabulazione}
    \label{rf:RFF-1}%
    \subreq{RFF-1.1}{\hyperref[uc:UC-9]{UC-9}}{Il sistema deve visualizzare un pulsante che permetta di iniziare l'analisi dell'ordine della squenza di tabulazione}
\end{freq}

\begin{freq}
    {RFF-2}
    {Il sistema deve poter visualizzare il risultato dell'analisi dell'ordine della sequenza di tabulazione}
    \label{rf:RFF-2}%

    \subreq{RFF-2.1}{\hyperref[uc:UC-9]{UC-9}}{Il sistema deve mostrare il numero totale di elementi tabbabili individuati}
    \subreq{RFF-2.2}{\hyperref[uc:UC-10]{UC-10}}{Il sistema deve mostrare una lista con le incongruenze individuate nell'ordine della sequenza di tabulazione}
\end{freq}

\subsubsection{Logica di business}

\paragraph{Obbligatori}

\begin{freq}
    {RFO-7}
    {Il sistema deve poter gestire la persistenza dei dati raccolti}
    \label{rf:RFO-7}%

    \subreq{RFO-7.1}{\hyperref[uc:UC-1]{UC-1}}{Il sistema deve poter salvare in modo persistente le informazioni relative alle immagini o ai video analizzati}
    \subreq{RFO-7.2}{\hyperref[uc:UC-1]{UC-1}, \hyperref[uc:UC-6]{UC-6}}{Il sistema deve poter ottenere le informazioni dei dati salvati}
\end{freq}

\begin{freq}
    {RFO-8}
    {Il sistema deve poter analizzare le immagini senza alt}
    \label{rf:RFO-8}%

    \subreq{RFO-8.1}{\hyperref[uc:UC-1]{UC-1}}{Il sistema deve poter ottenere una lista delle immagini presenti all'interno della pagina}
    \subreq{RFO-8.2}{\hyperref[uc:UC-1]{UC-1}}{Il sistema deve poter determinare la presenza dell'attributo \texttt{alt} per ogni immagine individuata}
    \subreq{RFO-8.3}{\hyperref[uc:UC-1]{UC-1}}{Il sistema deve poter ottenere il valore dell'attributo \texttt{src} delle immagini individuate}
    \subreq{RFO-8.4}{\hyperref[uc:UC-1]{UC-1}}{Il sistema deve poter calcolare il \texttt{base64} di un'immagine}
    \subreq{RFO-8.5}{\hyperref[uc:UC-1]{UC-1}, \hyperref[uc:UC-4]{UC-4}}{Il sistema deve poter fornire il \texttt{base64} di un'immagine all'interno della richiesta all'LLM}
    \subreq{RFO-8.6}{\hyperref[uc:UC-1]{UC-1}, \hyperref[uc:UC-2]{UC-2}}{Il sistema deve poter modificare lo stile CSS di un'immagine individuata}

\end{freq}

\begin{freq}
    {RFO-9}
    {Il sistema deve poter analizzare i video senza sottotitoli}
    \label{rf:RFO-9}%

    \subreq{RFO-9.1}{\hyperref[uc:UC-6]{UC-6}}{Il sistema deve poter ottenere una lista dei video presenti all'interno della pagina}
    \subreq{RFO-9.2}{\hyperref[uc:UC-6]{UC-6}}{Il sistema deve poter determinare la presenza di sottotitoli per ogni video individuato}
    \subreq{RFO-9.3}{\hyperref[uc:UC-6]{UC-6}}{Il sistema deve poter ottenere il valore dell'attributo \texttt{src} dei video individuati}
    \subreq{RFO-9.4}{\hyperref[uc:UC-6]{UC-6}}{Il sistema deve poter ottenere il binario di un video}
    \subreq{RFO-9.5}{\hyperref[uc:UC-6]{UC-6}, \hyperref[uc:UC-7]{UC-7}}{Il sistema deve poter fornire il binario di un video all'interno della richiesta all'LLM}
    \subreq{RFO-9.6}{\hyperref[uc:UC-6]{UC-6}, \hyperref[uc:UC-2]{UC-2}}{Il sistema deve poter modificare lo stile CSS di un video individuato}
\end{freq}

\begin{freq}
    {RFO-10}
    {Il sistema deve poter gestire errori nella connessione con l'LLM}
    \label{rf:RFO-10}%

    \subreq{RFO-10.1}{\hyperref[uc:UC-5]{UC-5}, \hyperref[uc:UC-8]{UC-8}}{In caso di indisponibilità dell'LLM, il sistema deve essere in grado di gestire l'errore in modo appropriato}
    \subreq{RFO-10.2}{\hyperref[uc:UC-8]{UC-8}}{Il sistema deve rilevare quando un errore è causato dall'invio di un video in un formato non supportato}
\end{freq}

\paragraph{Facoltativi}

\begin{freq}
    {RFF-3}
    {Il sistema deve poter analizzare l'ordine della sequenza di tabulazione}
    \label{rf:RFF-3}%

    \subreq{RFF-3.1}{\hyperref[uc:UC-9]{UC-9}}{Il sistema deve poter individuare tutti gli elementi tabbabili all'interno di una pagina}
    \subreq{RFF-3.2}{\hyperref[uc:UC-9]{UC-9}}{Il sistema deve poter ordinare gli elementi tabbabili secondo l'ordine della sequenza di tabulazione}
    \subreq{RFF-3.3}{\hyperref[uc:UC-9]{UC-9}, \hyperref[uc:UC-10]{UC-10}}{Il sistema deve poter ottenere il valore dell'attributo \texttt{tabindex} di un elemento della pagina, se presente}
    \subreq{RFF-3.4}{\hyperref[uc:UC-9]{UC-9}}{Il sistema deve poter determinare la posizione assoluta di un elemento all'interno della pagina}
    \subreq{RFF-3.5}{\hyperref[uc:UC-9]{UC-9}}{Il sistema deve essere in grado di identificare i contenitori semantici che racchiudono ciascun elemento}
    \subreq{RFF-3.6}{\hyperref[uc:UC-9]{UC-9}}{Il sistema deve poter determinare se due elementi tabbabili sono consecutivi nell'ordine della sequenza di tabulazione}
    \subreq{RFF-3.7}{\hyperref[uc:UC-9]{UC-9}}{Il sistema deve essere in grado di determinare la posizione relativa di un elemento rispetto a un altro, valutando se è più in alto o più a destra nella pagina}
    \subreq{RFF-3.8}{\hyperref[uc:UC-10]{UC-10}}{Il sistema deve poter modificare lo stile CSS di elemento tabbabile}
    \subreq{RFF-3.9}{\hyperref[uc:UC-10]{UC-10}}{Il sistema deve poter aggiungere nella pagina un badge per ogni elemento che indica la sua posizione nell'ordine della sequenza di tabulazione}

\end{freq}